\appendices
\section{Reference Experiments on Protocol Difficulty, Modality Sharing,
and Feature Physics}
\label{app:reference}

This appendix reports supporting experiments that bound the design decisions
of the main study. They deliberately use a generic window-level pipeline,
simpler than and independent of the constituent-level methodology of
Section~\ref{sec:methods}, so that the measured gaps reflect protocol
difficulty and representation structure rather than modeling choices. All of
them are released with the analysis code.

\subsection{Test Rig and Release Structure}
\label{app:rig}

Table~\ref{tab:rig} summarizes the machine, instrumentation, and release
organization underlying every experiment in this paper. Two properties drive
the protocol design. First, each diagnostic state is recorded once per
operating condition, so a run is the only sampling unit that is safely
independent; window-level splitting would place segments of one physical
recording on both sides of a split. Second, the twelve conditions are a
crossed grid of two excitation profiles, two load levels, and three speeds,
which makes it possible to hold out a condition either by interpolation or by
extrapolation---a distinction quantified in
Section~\ref{app:protocols}.

\begin{table}[!t]
\caption{Test rig, instrumentation, and release structure. Machine and sensor
parameters are as documented in the dataset descriptor \cite{chen2026}; the
release-structure rows are as observed in the distributed files.}
\label{tab:rig}
\centering
\footnotesize
\setlength{\tabcolsep}{4pt}
\begin{tabular}{l l}
\toprule
Property & Value \\
\midrule
Machine & QABP-90L2, 2.2~kW three-phase induction motor \\
Pole pairs & 2 \\
Single-side air gap & 0.5~mm \\
Bearings & SKF 6205-2Z-C3 (drive and non-drive end) \\
Bearing geometry & 9 balls, 7.94~mm ball, 39.04~mm pitch \\
Bearing orders & BPFO 3.585, BPFI 5.415, BSF 2.357 \\
Drivetrain & torque sensor, two-stage parallel gearbox, \\
           & magnetic powder brake \\
\midrule
Vibration & triaxial accelerometer, 100~mV/g \\
Current & three current clamps, 100~mV/A \\
Torque & S2001 sensor, 100~mV/Nm, $\pm$0.5\%~F.S. \\
Key-phase & one pulse per revolution \\
Sampling & 12.8~kHz, 8 synchronous channels, 90~s per run \\
\midrule
Diagnostic states & 24 (healthy, single, and compound) \\
Operating conditions & 12 (2 profiles $\times$ \{20, 40\}~Nm \\
                     & $\times$ \{1000, 2000, 3000\}~rpm) \\
Excitation profiles & speed circulation (constant load), \\
                    & torque circulation (constant speed) \\
Runs in the release & 288 (one per state and condition) \\
Compound runs & 108 over 9 compositions \\
\bottomrule
\end{tabular}
\end{table}

Three release-level properties were observed while preparing the data and are
reported because they materially affect protocol construction. (i) The
distributed archive contains 288 runs, whereas the descriptor article
reports 282; the surplus is a second \texttt{bearing\_outer\_H} recording at
the constant-torque 20~Nm, 3000~rpm condition, time-stamped roughly six weeks
after the main acquisition campaign. (ii) Compound-fault directory names
abbreviate the second constituent, so that a directory named
\texttt{bearing\_outer\_H\_and\_inner\_H} denotes an outer-race and
inner-race combination; parsing the labels naively splits one compound into
inconsistent constituent sets. (iii) Recordings of the same fault state made
on different dates are separable with near-perfect accuracy from the same
features used for diagnosis (run-level AUC 0.998--0.999). Property~(iii) is
the reason the transfer analyses in the main study are bounded by
recording-sensitivity and same-date controls rather than being read directly
as fault physics.

\subsection{Protocol Difficulty Under Operating-Condition Shift}
\label{app:protocols}

Windows of 0.64~s (8192 samples, non-overlapping) were described by 108 time-
and frequency-domain features (eight temporal statistics and ten spectral
descriptors per channel over the three vibration and three current channels)
and classified into the 24 diagnostic states by a 300-tree random forest, with
three seeds per configuration. Only the split changes across the six
protocols compared in Table~\ref{tab:condref}: \emph{random windows} splits
windows irrespective of recording (a deliberately leakage-prone reference);
\emph{in-condition} trains on the first 60\% of every recording and tests on
the final 30\% with a 10\% guard gap; \emph{unknown condition} holds out each
of the twelve operating conditions in turn; \emph{cross-profile} trains on one
excitation profile and tests on the other; \emph{steady~$\rightarrow$
transitional} trains on quasi-stationary segments and tests on speed and load
ramps, with stationarity labeled from the key-phase-derived speed and the
torque channel (used only for constructing this protocol, never as a
classifier input); and \emph{single source} inverts the unknown-condition
protocol, training on one condition and testing on the remaining eleven.

\begin{table}[!t]
\caption{Reference protocols, 24-class accuracy (mean $\pm$ standard deviation
over folds and three seeds). Features and classifier are identical
throughout; only the split changes.}
\label{tab:condref}
\centering
\footnotesize
\setlength{\tabcolsep}{5pt}
\begin{tabular}{l c c}
\toprule
Protocol & Folds & Accuracy \\
\midrule
Random windows (leakage-prone reference)   & 1  & $0.975 \pm 0.000$ \\
In-condition (guard-gap temporal split)    & 1  & $0.965 \pm 0.000$ \\
Unknown condition (leave-one-out)          & 12 & $0.939 \pm 0.024$ \\
Cross-profile                              & 2  & $0.839 \pm 0.115$ \\
Steady $\rightarrow$ transitional segments & 1  & $0.766 \pm 0.001$ \\
Single source (train on one condition)     & 12 & $0.333 \pm 0.104$ \\
\bottomrule
\end{tabular}
\end{table}

The same features and classifier span 0.975 to 0.333 across protocols.
Holding one condition out of twelve costs only ${\sim}2.6$ accuracy points
relative to the in-condition reference, because the eleven retained
conditions bracket the held-out one in speed and load and the model
interpolates. Training on a single condition and extrapolating to the other
eleven costs over 60 points; paired over the twelve condition folds, the
difference between the two directions is 0.606 (paired $t$-test
$p = 7.7\times10^{-10}$; Wilcoxon $p = 4.9\times10^{-4}$). Transfer between
excitation profiles is likewise asymmetric \emph{at equal training-set size}
(144 runs each): training under variable speed generalizes to variable load
(accuracy 0.944), whereas the reverse loses ${\sim}21$ points (0.734; paired
over seeds, $p = 3.3\times10^{-5}$). Because this contrast is size matched, it
isolates excitation diversity rather than data volume, unlike the
single-source contrast, which reduces both. Finally, appending 36
speed-invariant envelope-order descriptors leaves the interpolating protocol
unchanged ($+0.001$ on unknown condition, paired $p = 0.84$) but improves the
extrapolating one ($+0.047$ on single source, paired
$p = 8.4\times10^{-4}$), consistent with their construction.

Two implications carry into the main study. Methodologically, robustness to
operating conditions should be claimed in the extrapolating direction, which
is why the crossed design of Section~\ref{subsec:crossed_protocol} evaluates
composition and compound-context regimes jointly with unseen conditions rather
than relying on leave-one-condition-out alone. Practically, excitation
diversity during commissioning matters at least as much as data volume.

\subsection{Why Constituent Detectors Use Mechanism-Specific Features}
\label{app:modality}

The main study restricts each constituent detector to a curated,
mechanism-appropriate feature family (Section~\ref{subsec:physics_features})
rather than exposing every detector to the full multimodal bank. The
following experiment measures what a shared feature space does in this
dataset. Multi-label one-vs-rest random-forest detectors were trained on
single-fault windows and applied to all compound windows under three feature
views: vibration only, current only, and the concatenation of both.

The modality dissociation was complete: vibration features never detected an
electrically expressed constituent (detection rate 0.000 across all eight
cross-modal compounds) and current features never detected a mechanical one.
Critically, concatenation did not combine the two abilities but destroyed one
of them. The winding constituent was detected on 39.3\% of compound windows
from current features alone and on 0.2\% once vibration features were
concatenated---a roughly 200-fold collapse---while the mechanical constituent
retained 0.072 detection in the shared view. The electrical evidence is
therefore present and learnable in its own modality, and is lost to the
dominant mechanical signature when both families compete inside one
representation. Restricting each detector to physically appropriate
descriptors, as the main study does, removes that competition by
construction.

\subsection{Empirical Checks on Two Feature-Construction Choices}
\label{app:featurephysics}

Two parameters of the order-domain descriptors deserve explicit empirical
support rather than appeal to convention.

\emph{Integration bandwidth.} Order-domain descriptors integrate spectral
energy within $\pm0.12$ order of a target order. A fixed absolute window is a
constant \emph{relative} window only at the fundamental: at the second
harmonic it spans half the relative width, while the physical deviation of a
bearing order---rolling-element slip makes the true defect order differ from
its kinematic value---scales with the harmonic. To test whether this matters
in practice, the normalized band ratio was recomputed at half-widths of 0.06,
0.12, 0.25, and 0.50 order and scored by its ability to separate the
corresponding single-fault runs from healthy runs (run-level AUROC, median
over windows, 84 runs, order resolution 0.017).

The measured peak offsets confirm the concern and bound it. The local maximum
of the envelope-order spectrum lies a median 0.024 order from the kinematic
BPFO and 0.036 order from BPFI at the fundamental, but 0.049 and 0.080 order
at the second harmonic---the offset roughly doubles with harmonic number, as
slip-induced deviation should. Consequently the $\pm0.12$ window contains the
observed peak in 100\% and 96\% of runs for BPFO (first and second harmonic)
but in only 96\% and 74\% of runs for BPFI. The second-harmonic BPFI
descriptor is therefore the least reliable member of the bearing family.

Widening the window, however, does not repair this and measurably harms
discrimination. At the fundamental, BPFO separates faulty from healthy runs
perfectly (AUROC 1.000) at every half-width, so the choice is immaterial
there. At the second harmonic, AUROC for BPFO is 0.877 at $\pm0.12$ but falls
to 0.573 at $\pm0.25$ and to chance (0.497) at $\pm0.50$. The reason is
structural: $2\times$BPFO${}=7.17$ and $2\times$BPFI${}=10.83$ sit close to
integer shaft orders, which carry rotational and imbalance energy in healthy
and faulty machines alike, so a wider window admits common-mode energy faster
than it recovers the displaced fault peak. The $\pm0.12$ choice is thus close
to the best available compromise for this rig rather than an arbitrary one,
with the caveat that second-harmonic descriptors---BPFI's in
particular---carry less reliable evidence than their fundamentals. A
harmonic-scaled window would be the principled alternative, and is left to
future work because widening at fixed center is demonstrably not the fix.

\emph{Broken-bar sideband placement.} The current-domain descriptors are
evaluated at $o_e \pm 1$ and $o_e \pm 2$, that is, at offsets of one and two
\emph{shaft} orders from the electrical carrier. The classical broken-bar
signature instead lies at $(1 \pm 2s) f_e$, an offset of $2 s\, o_e$ shaft
orders, where the slip $s$ follows from the measured carrier and the
key-phase speed as $s = 1 - p/o_e$.

The pole count was determined from the data rather than assumed: the detected
carrier sits at $50.06$~Hz against a $49.69$~Hz shaft at the 3000~rpm setting
(and at $16.63$~Hz against $16.32$~Hz at 1000~rpm), so the machine is
two-pole, $p=1$, and the recovered slip is 3.0\%, 1.5\%, and 1.0\% at the
1000, 2000, and 3000~rpm settings---a near-constant slip \emph{speed} of about
$0.3$~Hz, as a roughly constant load torque implies.

The consequence is quantitative. The broken-bar sidebands lie $2 s\, o_e =
0.062$, $0.031$, and $0.020$ order from the carrier at the three speed
settings (median $0.024$ over 245 analyzed windows), which is only about two
spectral bins at the 2-s window length used here, and falls \emph{inside} the
$\pm0.12$ carrier band in 92.7\% of windows. The classical broken-bar
signature is therefore absorbed into the carrier descriptor itself rather
than resolved as a separate feature. The $o_e \pm 1$ and $o_e \pm 2$
descriptors, by contrast, sit one and two shaft orders from the carrier,
which is 16 to 50 times further out than the true sideband offset and
coincides instead with the $f_e \pm f_r$ spacing associated with
eccentricity. These descriptors should accordingly be read as
carrier-relative modulation proxies---which is how
Section~\ref{subsec:physics_features} interprets them---and not as
slip-parameterized broken-bar measurements. Resolving genuine broken-bar
sidebands on this rig would require either substantially longer analysis
windows or explicit slip-tracked demodulation, and remains open.
